\documentclass{beamer}
\usetheme{Madrid} % You can change this theme
\usepackage{graphicx}

\title{Implementation of Dielectric ELC in ESPResSo}
\subtitle{Bachelor's Thesis}
\author{Julien Hemminger}
\institute{Simulation Technology B.Sc.}
\date{April 2026}

\begin{document}

\begin{frame}
    \titlepage
\end{frame}

\begin{frame}{Motivation}
    There already is an ELC implementation that can handle dielectric interfaces in ESPResSo, but:
    \begin{itemize}
        \item It has errors that are poorly understood.
        \item The code is unmaintainable and very difficult to read.
    \end{itemize}
    \textbf{Goal:} Replace the old ELC implementation with a clean, fixed, modern version.
\end{frame}

\begin{frame}{Introduction}
\begin{figure}[htbp]
    \centering
    % Left Image
    \begin{minipage}{0.48\textwidth}
        \centering
        \includegraphics[width=\linewidth]{sandwich.png}
        
        \label{fig:sandwich}
    \end{minipage}
    \hfill % Adds horizontal space between the two images
    % Right Image
    \begin{minipage}{0.48\textwidth}
        \centering
        \includegraphics[width=\linewidth]{schematic_fully_periodic_replicated_slab_system.png}
        
        \label{fig:schematic}
    \end{minipage}
\end{figure}


    Developing in Python allows for faster iteration:
    \begin{itemize}
        \item Easier to develop, test, and prototype.
        \item Simplified visualization and validation of results.
    \end{itemize}
\end{frame}

\subsection{Regular ELC}

\begin{frame}{Regular ELC: Development Strategy}
    Iterative/test-driven development by breaking the problem down:
    \begin{itemize}
        \item \textbf{Workflow per system:}
        \begin{itemize}
            \item Compute reference solutions (analytical or brute-force).
            \item Adjust ELC prototype to match reference solutions.
        \end{itemize}
        \item \textbf{Systems analyzed:}
        \begin{itemize}
            \item Large boxes (negligible PBC influence with $L_x, L_y >> 1$).
            \item Smaller, neutral boxes (no non-neutral correction term yet).
            \item Full parameter sweeps (Non-neutral, Madelung, random, and extreme parameters).
        \end{itemize}
    \end{itemize}

    \vfill % Pushes the formula and footnote to the bottom

    $$E_{2D+h} = E_{3D} + E_{\text{corr}} + E_{\text{non-neutral}} + E_{\text{dipole}} \footnote{Tyagi, S., Arnold, A., \& Holm, C. (2008). J. Chem. Phys. 129, 204102.}$$
\end{frame}

\begin{frame}{Regular ELC: Results}
    \begin{figure}[h!]
        \centering
        \includegraphics[width=0.7\linewidth]{mar3_regular_elc_done.png}
        \caption{Validation of Regular ELC}
    \end{figure}
\end{frame}

% First Slide: Single Plate Status & Main Results
\begin{frame}{Dielectric ELC: Single Plate (Current Status)}
    \begin{itemize}
        \item Reference used: Ewald sum.
        \item Works for most parameters, though some outliers remain.
    \end{itemize}
    
    \begin{figure}
        \centering
        \includegraphics[width=0.7\linewidth]{apr27_single_plate_elcic_mostly_works_but_some_outliers.png}
        \caption{Single Plate ELCIC Results}
    \end{figure}
\end{frame}

% Second Slide: Neighbourhood Scan
\begin{frame}{Dielectric ELC: Neighbourhood Scan}
    \begin{figure}
        \centering
        \includegraphics[width=0.8\linewidth]{neighbourhood scan.png}
        \caption{Detailed Neighbourhood Scan Analysis}
        \label{fig:placeholder}
    \end{figure}
\end{frame}

\begin{frame}{Dielectric ELC: Dual Plates}
    \begin{itemize}
        \item Reference: Brute Force summation failed to converge.
    \end{itemize}
    
    \begin{figure}[h!]
        \centering
        \includegraphics[width=0.7\linewidth]{apr9_cant_get_dual_plate_ana_to_converge.png}
    \end{figure}
\end{frame}

\section{Implementing in ESPResSo}
\begin{frame}{Implementing in ESPResSo}
    \begin{itemize}
        \item Upcoming: C++ implementation.
        \item Integration with \textbf{Kokkos} support for hardware acceleration.
    \end{itemize}
\end{frame}

\end{document}